\chapter{Contexto do Projeto}

O planejamento do software foi estritamente atrelado ao contexto do negócio em que foi baseado, o que acabou também sendo a razão de apenas recursos essenciais serem adotados. Portanto, é importante tomar nota das características específicas nas quais o parceiro envolvido se encontra.

\section{Cenário}
A aplicação que será resultado do projeto será articulada no contexto operacional da organização, levando em consideração que ela é projetada para as suas atividades diárias. Assim, o proposto envolverá atividades como realização de vendas, manutenção do estoque com registro de aquisições, trocas ou até um inventário mais prático.

\section{Público-Alvo}
O presente projeto tem como proposta o desenvolvimento de uma solução computacional voltada à digitalização e gestão de acervo, tendo como parceiro o Sebo do Cesar - loja que comercializa livros usados, localizado no bairro da Penha, em São Paulo – SP.
O projeto foi desenvolvido em colaboração direta, permitindo que as necessidades reais do estabelecimento sejam consideradas durante a definição e evolução da solução. Dessa maneira, buscou-se desenvolver não apenas uma aplicação funcional, mas uma ferramenta alinhada ao contexto e aos processos efetivamente utilizados pelo estabelecimento.
O sistema será exclusivamente desenvolvido para uso único da empresa, o que significa que apenas a equipe do sebo terá acesso à ele mesmo se tratando de uma aplicação Web. O perfil destas pessoas geralmente dispensa designs especiais para pessoas não acostumadas ao meio digital.

\section{Necessidades Identificadas}
A iniciativa parte da identificação da necessidade de digitalização e organização do acervo físico atualmente existente no estabelecimento parceiro. A sua ausência pode dificultar processos relacionados à catalogação, localização, controle de estoque e comercialização dos itens disponíveis.
Dessa forma, pretende-se desenvolver uma plataforma capaz de apoiar a transformação desses processos, centralizando as informações do acervo e proporcionando maior organização e eficiência às atividades do estabelecimento. A solução poderá contemplar, conforme o levantamento de requisitos e as necessidades identificadas junto ao parceiro, funcionalidades relacionadas à gestão de vendas e consulta às informações do acervo.

\section{Regras de Negócio}
A solução será desenvolvida considerando o porte e as necessidades atuais da organização parceira, priorizando funcionalidades essenciais para a gestão do negócio e evitando a implementação de recursos que não agreguem valor ao contexto operacional do estabelecimento. Como diretriz de desenvolvimento, o sistema deverá apresentar uma interface simples, intuitiva e de fácil utilização, permitindo que as operações rotineiras sejam executadas de forma rápida e eficiente. Também deverão ser adotadas práticas que garantam a integridade, confiabilidade e segurança das informações armazenadas.
Além disso, o tratamento de dados pessoais dos clientes deverá respeitar a legislação vigente, especialmente os princípios estabelecidos pela Lei Geral de Proteção de Dados - LGPD \cite{brasil2018lgpd}, garantindo o uso adequado das informações coletadas pelo sistema. Por fim, todas as funcionalidades implementadas deverão estar alinhadas aos requisitos levantados junto ao parceiro, mantendo o foco na digitalização do acervo, no controle das operações e no apoio à gestão do Sebo do Cesar.